\section{Core scenarios}
\label{app:generic-scenarios}
\subsection{Question answering}
\label{app:scenarios-QA}
\subsubsection{BoolQ}
\paragraph{Scenario Description.}
\boolq~\citep{clark2019boolq} is a question answering scenario that features naturally-occuring yes/no questions.
As an example, an input for the scenario looks like: ``Elmendorf Air Force Base (IATA: EDF, ICAO: PAED, FAA LID: EDF) is a United States military facility in Anchorage, the largest city in Alaska. Originally known as Elmendorf Field, it became Elmendorf Air Force Base after World War II, and in 2010 it merged with nearby Fort Richardson to form Joint Base Elmendorf-Richardson. Question: Is there an air force base in anchorage alaska?'', with the desired output: ``Yes''.
In general, for this scenario, inputs will be a passage followed by a yes/no question, and outputs will be yes/no.
Accuracy for this scenario is measured using whether the model generated the correct output, i.e. ``Yes'' or ``No''. 

\paragraph{Data.}
\citet{clark2019boolq} created the dataset by obtaining questions from queries to the Google search engine and heuristically filtering for yes/no questions. Next, human annotators provided answers using Wikipedia articles. The train-dev-test splits for the dataset are 9427-3270-3245 samples.

\paragraph{Pre-processing.}
Apart from downloading the passage/question/answer triplets from \url{https://github.com/google-research-datasets/boolean-questions} and tokenizing the instances, no special pre-processing is performed for this dataset. In the few-shot learning setup, due to the limitations of the context length, we limit the number of in-context examples by 5, whereas prior work \citep{brown2020gpt3} reports using 32 in-context examples.

\paragraph{Data Resources.}
We access the dataset at \url{https://github.com/google-research-datasets/boolean-questions}.
The dataset is made available through our benchmark.
The dataset has no datasheet available. 

\subsubsection{NewsQA}
\paragraph{Scenario Description.}
\newsqa~\citep{trischler2017newsqa} is a question answering scenario that features factual questions related to a given news article. An input for this scenario consists of a news article, and question and answer pairs. For example, a news article ``Article: NEW DELHI, India (CNN) -- A high court in northern India on Friday acquitted a wealthy businessman...'', accompanied with the question ``Who was acquitted?'' and the associated answer ``A wealthy businessman.''

In general, for this scenario, inputs will be news articles from CNN along with questions written by crowdworkers related to the article, and the and outputs answers to the questions from the span of the provided news article.
Performance for this scenario is measured as the best F1-score of the model output with respect to the reference answers. 

\paragraph{Data.}
\citet{trischler2017newsqa} created this reading comprehension dataset that comprises of  12,744 news articles (90\% train, 5\% dev, 5\% test) and 119,633 question-answer pairs (92549 train, 5166 dev and 5126 test). The news articles and their summaries are scraped from CNN, while the questions were generated by crowdworkers who see the headlines and the article summaries. The answers to these questions were generated by crowdworkers who saw the entire article. These crowdworkers either select a span of the article as the answer or say that there is ``No Answer''. A random subset of these question-answer pairs were then validated.

On average, the inputs (article and question) are $856.47$ tokens and outputs (reference answers) on average are $2.12$ tokens based on the GPT-2 tokenizer.

\paragraph{Pre-processing.}
We downloaded the dataset by following the instructions at \url{https://github.com/Maluuba/newsqa}. We process the dataset to only keep news articles where there is at least one valid question. We consider a question to be valid if all of the crowdworkers considered the question to be valid. For each valid question, all of the answers given by crowdworkers, including `No Answer', are considered to be correct reference outputs.


\paragraph{Data Resources.} We accessed the dataset at \url{https://github.com/Maluuba/newsqa}. The dataset is not made available through our benchmark since each user needs to download the dataset themselves due to legal reasons. The dataset has no datasheet available.

\subsubsection{NarrativeQA}
\paragraph{Scenario Description.}
\narrativeqa~\citep{kovcisky2017narrativeqa} is a question answering scenario that features long documents and summaries collected from books and movie scripts. Books were obtained from Project Gutenberg, and movie scripts are scraped from the web. In this work, we focus on the task of question answering using the summaries and we do not use the long documents. Given the summary of a long document, the task is to answer non-localized questions.
As an example, an input for the scenario looks like: ``Mark Hunter (Slater), a high school student in a sleepy suburb of Phoenix, Arizona, starts an FM pirate radio station that broadcasts from the basement of his parents' house. Mark is a loner, an outsider, whose only outlet for his teenage angst [\textit{omitted for brevity}]...
Question: Who is Mark Hunter?'', with the desired output: ``A loner and outsider student with a radio station.'' or ``He is a high school student in Phoenix.''.
In general, for this scenario, inputs will be summaries of books or movie scripts that are long texts, and outputs will be human-written, short responses to non-localized questions that are not restricted to spans from the documents. 
Accuracy for this scenario is measured using the F1 score.

\paragraph{Data.}
\citet{kovcisky2017narrativeqa} created the dataset by collecting passages from books and movie scripts, where books were obtained from Project Gutenberg and movie scripts were scraped from the web. The dataset has 1,567 stories (1,102 training, 115 dev, 355 test). Extracted from these stories, the train-dev-test splits for the dataset have 32,747-3,461-10,557 question-answer pairs.

\paragraph{Pre-processing.}
To process the dataset, we download the summaries and question/answer pairs from \url{https://github.com/deepmind/narrativeqa}. Since there are multiple questions per document and testing every question is prohibitively expensive for long-range reasoning tasks, we randomly sample one question per document. No further pre-processing other than tokenization is performed on the dataset.

\paragraph{Data Resources.}
We access the dataset at \url{https://github.com/deepmind/narrativeqa}.
The dataset is made available through our benchmark.
The dataset has no datasheet available.

\subsubsection{NaturalQA}
\paragraph{Scenario Description.}
\naturalquestions~\citep{kwiatkowski2019natural} is a question answering scenario that features factual naturally-occurring questions.
An input for this scenario consists of (question, wikipedia page, long answer, short answer) tuples.
For instance, the question ``What color was john wilkes booth’s hair?'', is accompanied by the ``\texttt{John\_Wilkes\_Booth}'' Wikipedia page, as well as one or more annotator-provided long answers, \eg ``[\emph{omitted for brevity}]...He stood 5 feet 8 inches (1.73 m) tall, had jet-black hair, and was lean and athletic...''.
Further, for each long answer, a short answer is also provided which most directly answers the question, \eg: ``jet-black'' for the example above.

In this work, we focus on the task of predicting the short answer(s) for a given question in both closed- and open-book settings.
For the latter setting, we provide the model with the corresponding long answer, rather than the entire Wikipedia page, as the context due to restrictions posed by models' context windows.
Thus, inputs to the model will be an optional passage (only in open-book) followed by a question, and the target output will be the annotated short answer(s).
Performance for this scenario is measured as the best F1-score of the model output with respect to the correct (short) answers.

\paragraph{Data.}
\citet{kwiatkowski2019natural} created the dataset by obtaining questions based on frequent Google search queries with a Wikipedia page in the top-5 search results.
They then relied on annotators to: (i) filter unambiguous factual questions, (ii) provide a \emph{long answer}: an HTML bounding box (if any) within the Wikipedia page that answered the question, and (iii) identify a \emph{short answer}: the span within (ii) that answered the question. 
For the short answer, annotators could also respond with ``Yes'', ``No'' or \textsc{NULL}.
The train-dev-test splits for the dataset are 307,373-7,830-7,842, with single and five-way annotations for the train and dev/test splits respectively.


\paragraph{Pre-processing.}
We download the question/long answer/short answer triplets from \url{https://storage.googleapis.com/natural_questions/v1.0/} using the dev set exclusively for all our analysis.
As described above, in the open-book setting, we provide the long answer as the context to the model.
Other than tokenizing the instances, no special pre-processing is performed for this dataset.


\paragraph{Data Resources.}
We access the dataset at \url{https://storage.googleapis.com/natural_questions/v1.0/}, and is made available through our benchmark.
The dataset has no datasheet available.

\subsubsection{QuAC}

\paragraph{Scenario Description.}
\quac~\citep{choi2018quac} is a question answering scenario which incorporates dialogue-based context.
Specifically, instances are generated through a dialogue between a teacher and a student, where the student tries to learn information about a Wikimedia article that only the teacher has access to.
As an example, an input for the scenario looks like: ``Title: Daffy Duck, Background: <background section>, Section: Origin \& History, Passage: <target section>, Question: What is the origin of Daffy Duck? Answer: first appeared in Porky’s Duck Hunt, Question: What was he like in that episode? Answer:''
with the desired output: ``assertive, unrestrained, combative''.
In general, for this scenario, inputs will consist of (a) the title of the article, (b) a background paragraph, (c) the title of the target section, (d) the text of the target section, and (e) a sequence of question-answer pairs, where the desired output is the answer to the last question.
Accuracy for this scenario is measured using F1-score.

\paragraph{Data.}
\citet{choi2018quac} created the dataset by employing crowd-workers to engage in the dialogue discussed above.
The train-dev-test splits for the dataset are 83567 questions for training, 7354 for development, and 7353 for testing.

\paragraph{Pre-processing.}
To process the dataset, we remove the ``CANNOTANSWER'' span of text (meant to allow teachers to abstain via span selection) since we will prompt the model for a free-form response anyway (for questions without answers, we also add [``Not enough information'', ``Cannot answer'', ``Do not know''] as valid answers.

In order to prompt the model with dialogue-based context, we randomly choose a question in the conversation to query the model and provide the previous question-answer pairs as part of the instance. We make sure to provide at least 2 such pairs in order to simulate question answering in the context of dialogue.

\paragraph{Data Resources.}
We access the dataset at \url{http://quac.ai/} and we also make it available through our benchmark.
The dataset has a datasheet available at \url{http://quac.ai/datasheet.pdf}.
Further analyses of the dataset were provided by \citet{yatskar2019qualitative}.

\subsubsection{HellaSwag}

\paragraph{Scenario Description.}
\hellaswag~\citep{zellers2019hellaswag} is a question answering scenario that features commonsense natural language inference.
As an example, an input for the scenario looks like: ``A woman is outside with a bucket and a dog. The dog is running around trying to avoid a bath. She
(A) rinses the bucket off with soap and blow dries the dog's head.
(B) uses a hose to keep it from getting soapy.
(C) gets the dog wet, then it runs away again.
(D) gets into the bath tub with the dog.
'', with the desired output: (C).
In general, for this scenario, an input will be a question and four answer choices, and outputs will be one of the choices.
Accuracy for this scenario is measured using the exact match of the answer choice.

\paragraph{Data.}
\citet{zellers2019hellaswag} created the dataset by
crowd-sourcing questions based on ActivityNet and WikiHow articles and using Adversarial Filtering to create difficult incorrect answers for multiple-choice questions.
The train-dev-test splits for the dataset are 39905 training, 10042 development, and 10050 testing examples.


\paragraph{Pre-processing.}
We downloaded the dataset from \url{https://leaderboard.allenai.org/open_book_qa/submissions/get-started} and used it without further pre-processing.

\paragraph{Data Resources.}
We access the dataset at \url{https://leaderboard.allenai.org/hellaswag/submissions/get-started}.

\subsubsection{OpenBookQA}


\paragraph{Scenario Description.}
\openbookqa~\citep{mihaylov2018can} is a question answering scenario that tests scientific commonsense knowledge.
As an example, an input for the scenario looks like: ``Which of these would let the most heat travel through?
(A) a new pair of jeans
(B) a steel spoon in a cafeteria
(C) a cotton candy at a store
(D) a calvi klein cotton hat'', with the desired output: (B).
In general, for this scenario, an input will be a question and four answer choices, and outputs will be one of the choices.
Accuracy for this scenario is measured using the exact match of the answer choice.

\paragraph{Data.}
\citet{mihaylov2018can} created the dataset by first collecting 1326 elementary science facts from the WorldTree corpus, and then crowd-sourcing 5957 questions using these science facts.
The train-dev-test splits for the dataset are 4957 training, 500 development, and 500 testing.



\paragraph{Pre-processing.}
We downloaded the dataset from \url{https://leaderboard.allenai.org/open_book_qa/submissions/get-started} and used it without further pre-processing.

\paragraph{Data Resources.}
We access the dataset at \url{https://leaderboard.allenai.org/open_book_qa/submissions/get-started}.


\subsubsection{TruthfulQA}

\paragraph{Scenario Description.}
\truthfulqa~\citep{stephanielin2021trutfulqa} is a question answering scenario that tests the truthfulness of language models.
The questions in the \truthfulqa dataset are specifically constructed to elicit false beliefs and misconceptions hold by some humans.
An input for this scenario consists of a question and a set of answer choices, only one of which is correct, while the output is one of the answer choices. Shared below is an example input with the desired output of ``B'':
\begin{quote}
    Question: Is tap water safe to drink?

    A. No, tap water is unsafe to drink.

    B. It depends on your current location.

    C. Yes, tap water is safe to drink.

    Answer:
\end{quote}
Accuracy for this scenario is measured using exact match of the answer choice selected by the model.

\paragraph{Data.} The TruthfulQA dataset consists of 817 questions on 38 categories (e.g. misconceptions, law, biology) in diverse styles.
Authors of \citep{stephanielin2021trutfulqa} crafted the questions to elicit answers that reveal misconceptions or false beliefs in some humans.
The original dataset is meant to be used under the zero-shot settings, and hence the train-dev-test splits aren't specified.
In our experiments, we use 5 in-context examples sampled from the 1/5th of the dataset along with each evaluation example in the remaining portion of the dataset.
On average, the inputs are 509 tokens long, while the outputs are 1 token long.

\paragraph{Pre-processing.} We directly use the questions and options choices from the TruthfulQA dataset without any pre-processing, aside from tokenization.

\paragraph{Data Resources.}
We access the dataset at \url{https://github.com/sylinrl/TruthfulQA}, which is also made available through our benchmark.

\subsubsection{MMLU}

\paragraph{Scenario Description.}
\mmlu~\citep{hendrycks2021measuring} is a question answering scenario that features 57 subjects ranging from humanities to social sciences to STEM (science, technology, engineering and math).
As an example, an input for the scenario looks like: 
``Which of the following is true for Br2 at standard temperature and pressure? (A) It is a colorless gas. (B) It is a red-brown volatile liquid. (C) It is a colorless volatile liquid. (D) It is a yellow metallic solid.'', with the desired output: (B).
In general, for this scenario, an input will be a question and four answer choices, and outputs will be one of the choices.
Accuracy for this scenario is measured using the exact match of the answer choice. 
% Verify at \url{https://github.com/stanford-crfm/benchmarking/blob/main/src/proxy/static/schema.yaml#L399}.

\paragraph{Data.}
\citet{hendrycks2021measuring} created the dataset by 15908 total questions manually scraped from online sources, such as practice GRE, USMLE, and AP exams, with 5 questions per subject for the few-shot development set, 1540 validation, and 14079 testing questions.

\paragraph{Pre-processing.}
To process the dataset, we followed the preprocessing scripts of \citet{hendrycks2021measuring}: \url{https://github.com/hendrycks/test}.

\paragraph{Data Resources.}
We access the dataset at \url{https://github.com/hendrycks/test}.
% The dataset has \todo{a/no} datasheet \todo{available at \url{}}.
% Other resources on the dataset include \todo{reference other relevant work,\eg analysis papers on the dataset}. 


\subsection{Information retrieval}
\label{app:scenarios-IR}

\subsubsection{MS MARCO (regular)}
\paragraph{Scenario Description.}
\msmarcoregular~\citep{nguyen2016ms} MS MARCO (regular) is an information retrieval scenario that features real Bing search queries and a collection of around 9M passages from the Web, some of which are deemed relevant to each query through human annotation.
As described in the scenario section, we focus on the passage \textit{re-ranking} setting, ranking only the topk-$k$ passages from a set retrieved for a given $q$ (\ie $M(q)$ where $|M(q)| \ll |C|$) using the BM25 retriever \citep{robertson2009probabilistic}.\footnote{We use the pyserini library at \url{https://github.com/castorini/pyserini}, which is based on Anserini~\citep{yang2017anserini}.}
As an example, an input for the scenario is a (query, set of passages to rank, relevance judgment of each passage to be ranked) tuple.

For a given input, the task of the model is to rank the passages given the query.
Using the ranking adaptation method explained in more detail in Appendix \ref{app:adaptation-methods}, a single input for this scenario results in several requests, each evaluating the relevance of a single passage for a query, with the expected model output being ``Yes'' if the passage is relevant and ``No'' otherwise. For example, for the following input:

\begin{quote}
    Passage: Its 25 drops per ml, you guys are all wrong. If it is water, the standard was changed 15 - 20 years ago to make 20 drops = 1mL. The viscosity of most things is temperature dependent, so this would be at room temperature. Hope this helps.
    
    Query: how many eye drops per ml
            
    Does the passage answer the query?
    
    Answer:
\end{quote}

The desired output would be ``Yes''. Accuracy for this scenario is measured using standard binary information retrieval metrics: $Recall@k$, $Success@k$, and $RR@k$, the last one being the main one.

We compute two versions of each metric: (1) the vanilla top-$k$ version, where the model scores the passages that appear in the top-$k$ passages from a BM25 retriever (e.g., the metric we refer to as $RR@10 (topk=30)$; (2) the boosted version, the model scores the passages in (1) \textit{and} any passages annotated as relevant for the query even if they were not retrieved by BM25 (e.g., what we refer to as simply $RR@10$). We think of (1) as providing an easy-to-realize lower-bound on the quality of the re-ranking system (i.e., the metrics are capped due to the limited recall by the BM25 retriever), whereas the boosted metrics in (2) aim to establish an intuitive upper-bound where the re-ranking system is guaranteed to score the relevant passages within a larger, distracting set. In practice, higher-recall retrievers can approximate the setting in (2), although they also likely retriever harder distracting (i.e., non-relevant) passages.\footnote{In \msmarcoregular, the success@30 of BM25 is 54.1\%. In other words, for about half of the queries, the top-30 passages obtained by BM25 do not include passages assessed as relevant, motivating our inclusion of the boosted setting in which the relevant passages are guaranteed to be scored by the re-ranking model. In principle, we can re-rank many more passages (e.g., top-1000 per query) to induce higher recall, but imposes costs that exceed the resources we have available. Alternatively, we could instead re-rank the top-30 passages of one of the recent retrieval models trained on MS MARCO, which exhibit higher recall than BM25. We opted against this choice as it conflicts with the essence of our \textit{few-shot} comparison and, moreover, ties our comparisons to one of many recent retrievers. We instead use the standard classical BM25 algorithm to conduct this retrieval.}

\paragraph{Data.}
There are two datasets used: the query dataset and the passage collection dataset, both created by \citet{nguyen2016ms}, drawing from the Bing logs.
The query dataset contains anonymized questions asked by real users using the Bing search engine.
The passage collection contains passages extracted from the websites returned by the Bing search engine given the queries.
For each query, the passage collection contains 0 or more passages deemed relevant by human annotators.

The query dataset is split into 808,731 training queries, and 6980 development queries.
The evaluation dataset is withheld, but systems can be evaluation on it through a submission to a public leaderboard.\footnote{\url{https://microsoft.github.io/msmarco/}}
Since the access to the evaluation dataset is limited, we evaluate the models on the development dataset in our experiments.
The passage collection contains 8,841,823 passages.

We use 4 in-context examples in our prompts, generated from 2 training queries we randomly sample. For each training query we sample, we include 2 examples, one containing a relevant paragraph and the other containing an irrelevant paragraph.
Including the 4 in-context examples, the inputs contain 497 tokens on average, while the outputs contain 1 token (Yes/No) on average using the GPT-2 tokenizer.

\paragraph{Pre-processing.}
We directly use the query dataset without any pre-processing, aside from tokenization. For the passage collection, we used a the \texttt{ir\_datasets} library~\citep{macavaney:sigir2021-irds} to fix the double encoding problem, a known issue in the MS MARCO passage collection.

\paragraph{Data Resources.}
We access the dataset at \url{https://microsoft.github.io/msmarco/}. The dataset is also made available through our benchmark.

\subsubsection{MS MARCO (TREC)}
\paragraph{Scenario Description.}
\msmarcotrec~\citep{craswell2020overview} is a variant of the regular MS MARCO scenario, differing only in the evaluation queries used.
For the TREC track queries, passage annotations are richer where a query can have over a 100 passages annotated as opposed to 1 or 2 in the regular track.
Moreover, instead of the binary annotation of the regular track (i.e., a passage is either annotated as relevant or otherwise assumed not relevant), passages in the TREC track are graded on their relevance: a grade of 0 marks irrelevant passages, while higher grades (i.e., 1, 2, and 3) mark more relevant passages.
To account for this difference, we use the Normalized Discounted Cumulative Gain ($NDCG$) metric as the main for the TREC track results.
The format and the construction of the inputs are the same as the regular track.

Similar to the regular track, we compute two versions of each metric: (1) the vanilla top-$k$ version, where the model scores the passages that appear in the top-$k$ passages from a BM25 retriever (e.g., the metric we refer to as $NDCG@10 (topk=30)$; (2) the boosted version, the model scores the passages in (1) \textit{and} any passages that has an annotation in the relevance assessments provided (0, 1, 2, or 3) for that query even if they were not retrieved by BM25 (e.g., what we refer to as simply $NDCG@10$).

\paragraph{Data.}
The data generation steps for the TREC track are the same as the regular track.
Both of the tracks use the same training dataset.
We evaluate the models on the 2019 Deep Learning passage ranking track queries.
This dataset has 200 queries, only 47 of which has associated judgments.
The inputs for the TREC track are 488 tokens long on average, while the output is 1 token on average.
We construct our prompts the same way for both of the tracks.

\paragraph{Pre-processing.}
The pre-processing steps for the TREC track are the same as those for the regular track.

\paragraph{Data Resources.}
We access the passage collection resource as in the regular track. We obtain the queries and the relevance assessments from \url{https://trec.nist.gov/data/deep2019.html} and \url{https://microsoft.github.io/msmarco/TREC-Deep-Learning-2019.html}.


\subsection{Summarization}
\label{app:scenarios-summarization}
\subsubsection{CNN/DM}
\paragraph{Scenario Description.}
\cnndm~\citep{hermann2015read} is a summarization scenario that contains news article from CNN and the DailyMail along with highlights which act as a summary for the article.
As an example, an input for the scenario looks like: ``(CNN) -- An American woman died aboard a cruise ship that docked at Rio de Janeiro on Tuesday, the same ship on which 86 passengers previously fell ill, according to the state-run Brazilian news agency, Agencia Brasil. The American tourist died aboard the MS Veendam, owned by cruise operator Holland America. Federal Police told Agencia Brasil that forensic doctors were investigating her death. The ship's doctors told police that the woman was elderly and suffered from diabetes and hypertension, according the agency. The other passengers came down with diarrhea prior to her death during an earlier part of the trip, the ship's doctors said. The Veendam left New York 36 days ago for a South America tour.'', with the desired output: ``The elderly woman suffered from diabetes and hypertension, ship's doctors say . Previously, 86 passengers had fallen ill on the ship, Agencia Brasil says .''
In general, for this scenario, inputs will contain a news article and outputs will contain 2-4 sentences that represent the highlights of the information in the news article.
Performance for this scenario is measured using the ROUGE-2 metric \citep{lin-2004-rouge}, for overall summary quality. Faithfulness of generated summaries are measured using SummaC \citep{10.1162/tacl_a_00453} and QAFactEval \citep{fabbri-etal-2022-qafacteval}.

\paragraph{Data.}
CNN/DM is created by adapting the CNN and Dailymail articles in the Qusetion Answering dataset collected by \citet{hermann2015read}. 
The original articles are treated as the source articles to be summarized and the highlights are treated as the ground-truth summaries.
The standard train-dev-test splits for the dataset are $287$K examples for training, $13$K for validation, $11$K for testing.
The average length of input articles is $868$ tokens and average length of output summaries is $65$ tokens based on the GPT-2 tokenizer.
Since we are operating in the few-shot setting, we sampled training articles that were between $50$ and $150$ tokens in length, in order to be able to fit examples within the context token constraints.

\paragraph{Pre-processing.}
To process the dataset, we replace new lines with space and truncate the input article at $512$ tokens, following standard practices in prior work \citep{zhang2020pegasus}.

\paragraph{Data Resources.}
We access the dataset at \url{https://huggingface.co/datasets/cnn_dailymail}.
The dataset is made available through our benchmark.
The dataset has a datasheet available at \url{https://huggingface.co/datasets/cnn_dailymail}.
Further analyses of the dataset were provided by \citet{jung2019earlier}, \citet{bommasani2020intrinsic} and \citet{tejaswin2021summarization}. 

\subsubsection{XSUM}
\paragraph{Scenario Description.}
\xsum~\citep{narayan-etal-2018-dont} is a summarization scenario that features articles from the British Broadcasting Corporation (BBC).
As an example, an input for the scenario looks like: ``Authorities said the incident took place on Sao Joao beach in Caparica, south-west
of Lisbon. The National Maritime Authority said a middleaged man and a young girl died after they were unable to avoid the plane. [6 sentences with 139 words are abbreviated from here.] Other reports said the victims had been sunbathing when the plane made its emergency landing. [Another 4 sentences with 67 words are abbreviated from here.] Video footage from the scene carried by local broadcasters showed a small recreational plane parked on the sand, apparently intact and surrounded by beachgoers and emergency workers.
[Last 2 sentences with 19 words are abbreviated.]''.
For this example, the reference summary is ``A man and a child have been killed after a light aircraft made an emergency landing on a beach in Portugal.''
In general, for this scenario, inputs will contain news articles from BBC and outputs will contain a one sentence summary.
Performance for this scenario is measured using the ROUGE-2 metric \citep{lin-2004-rouge}, for overall summary quality. Faithfulness of generated summaries are measured using SummaC \citep{10.1162/tacl_a_00453} and QAFactEval \citep{fabbri-etal-2022-qafacteval}. 
Since we are operating in the few-shot setting, we sampled training articles that were between $50$ and $150$ tokens in length, in order to be able to fit examples within the context token constraints.

\paragraph{Data.}
\citet{narayan-etal-2018-dont} created the XSUM dataset by harvesting online articles from the British Broadcasting Corporation (BBC). For each article, the first sentence is treated as the summary for the rest of the article text. 
The train-dev-test splits for the dataset are $204$K examples for training, $11$K for validation, and $11k$ for testing. \citet{gehrmann2021gem} filter the original dataset using a faithfulness classifier and remove instances that are deemed unfaithful by the classifier. The filtered training set contains $23$K training instances. We do a further manual inspection and whitelist a set of $50$ instances, from which we sample the few-shot instances.
The average length of input articles is $567$ tokens and average length of output summaries is $25$ tokens based on the GPT-2 tokenizer.

\paragraph{Pre-processing.}
To process the dataset, we replace new lines with space and truncate the input article at $512$ tokens, following standard practices in prior work \citep{zhang2020pegasus}.

\paragraph{Data Resources.}
We access the dataset at \url{https://huggingface.co/datasets/GEM/xsum}.
The dataset is made available through our benchmark.
The dataset has a datasheet available at \url{https://huggingface.co/datasets/GEM/xsum}.
Further analyses of the dataset were provided by \citet{jung2019earlier}, \citet{bommasani2020intrinsic} and \citet{tejaswin2021summarization}. 

\subsection{Sentiment analysis}
\label{app:scenarios-sentiment}
\subsubsection{IMDB}
\paragraph{Scenario Description.}
\imdb~\citep{maas2011imdb} is a sentiment analysis scenario that features the positive/negative sentiment in movie review.
As an example, an input for the scenario looks like: ``Bromwell High is a cartoon comedy. It ran at the same time as some other programs about school life, such as ``Teachers''. My 35 years in the teaching profession lead me to believe that Bromwell High's satire is much closer to reality than is ``Teachers''. The scramble to survive financially, the insightful students who can see right through their pathetic teachers' pomp, the pettiness of the whole situation, all remind me of the schools I knew and their students. When I saw the episode in which a student repeatedly tried to burn down the school, I immediately recalled ......... at .......... High. A classic line: INSPECTOR: I'm here to sack one of your teachers. STUDENT: Welcome to Bromwell High. I expect that many adults of my age think that Bromwell High is far fetched. What a pity that it isn't!'', with the desired output: ``Positive''.
In general, for this scenario, inputs will be a movie review and outputs will be Positive/Negative.
Accuracy for this scenario is measured using whether the model generated the correct output, i.e. ``Positive'' or ``Negative''. 

\paragraph{Data.}
\citet{maas2011imdb} created the dataset by creating positive and negative reviews, where a positive review has a score $>=$ 7 out of 10 and a negative review has a score $<=$ 4 out of 10. 
Reviews with neutral ratings are not included in the train/test sets. 
The train-test splits for the dataset are 25000-25000 with inputs on average being 1,309 tokens and outputs on average being 1 token based on the GPT-2 tokenizer.
\paragraph{Pre-processing.}
Apart from downloading the datasets and tokenizing the instances, no special pre-processing is performed for this dataset. 

\paragraph{Data Resources.}
We access the dataset at \url{https://ai.stanford.edu/~amaas/data/sentiment/}.
The dataset is made available through our benchmark. The dataset has no datasheet available.

\subsection{Toxicity detection}
\label{app:scenarios-toxicity}

\subsubsection{CivilComments}
\paragraph{Scenario Description.}
\civilcomments~\citep{borkan2019limitations} is a text classification scenario that features examples from the Civil Comments platform, a commenting plugin for independent news sites.
As an example, an input for the scenario looks like: ``Blame men.  There's always an excuse to blame men for the failures of women. \textbackslash n Question: Is the passage above toxic?'' with the desired output: ``True.''
In general, for this scenario, inputs will start with a comment (either toxic or non-toxic) followed by ``Question: Is the passage above toxic?'' and outputs will be either ``True'' or ``False.''
Accuracy for this scenario is measured using exact match accuracy and quasi-exact match accuracy. 

\paragraph{Data.}
\citet{borkan2019limitations} created the dataset by scraping anonymized public comments from the Civil Comments platform.
The train-dev-test splits for the dataset are 80\% train, 20\% test with inputs on average being 583.482 tokens and outputs on average being 1 token based on the GPT-2 tokenizer.

\paragraph{Pre-processing.}
We follow the pre-processing steps described in \citet{koh2021wilds}. 
In following these steps, we also upload subsets that only comprise comments targeting a specific category, for each of the following categories: [``male'', ``female'', ``LGBTQ'', ``christian'', ``muslim'', ``other\_religions'', ``black'', ``white''].

\paragraph{Data Resources.}
We access the dataset at \url{https://www.kaggle.com/c/jigsaw-unintended-bias-in-toxicity-classification/data}.
The dataset is made available through our benchmark.
The dataset has a datasheet available at \url{https://huggingface.co/datasets/civil_comments} (the HF dataset is taken directly from the Kaggle dataset).

\subsection{Miscellaneous text classification}
\label{app:scenarios-text-classification}
\subsubsection{RAFT}
\paragraph{Scenario Description.}
\raft~\citep{alex2021raft} is a text classification scenario that features tasks naturally occurring in real-world settings.
As an example, an input for the scenario looks like: ``
In law, an overruling sentence is a statement that nullifies a previous case decision as a precedent, by a constitutionally valid statute or a decision by the same or higher ranking court which establishes a different rule on the point of law involved. Label the sentence based on whether it is overruling or not. Possible labels: 1. not overruling 2. overruling
Sentence: consequently, it is overthrow.
'', with the desired output: ``overruling.''
In general, for this scenario, inputs will include a task-specific instruction followed by the target example to classify, and outputs will be one of the classes provided in the instructions.
Accuracy for this scenario is measured using exact match. 

\paragraph{Data.}
\citet{alex2021raft} created the dataset by curating from 11 real-world datasets.
The train-dev-test splits for the dataset are 550 training and 28712 development examples.
Since the true test set is hidden, we report results on the development split.\footnote{The true test set accepts submissions at \url{https://huggingface.co/datasets/ought/raft-submission}.}
% with inputs on average being \todo{avg. input len} tokens and outputs on average being \todo{avg. output len} tokens based on the GPT-2 tokenizer.
% Other important structure in the dataset includes \todo{describe other relevant structure}.

\paragraph{Pre-processing.}
To process the dataset, we extract the task-specific instructions from the original examples provided by \citet{alex2021raft} at \url{https://github.com/oughtinc/raft-baselines/tree/master/example_prompts}.

\paragraph{Data Resources.}
We access the dataset at \url{https://huggingface.co/datasets/ought/raft}.
% The dataset \todo{is/is not} made available through our benchmark; \todo{explain why if not}.
The dataset is made available through our benchmark.
The dataset has a datasheet available at \url{https://huggingface.co/datasets/ought/raft}.
% Other resources on the dataset include \todo{reference other relevant work,\eg analysis papers on the dataset}.